\documentclass[10pt,a4paper]{amsart}
\usepackage[margin=19mm]{geometry}
\usepackage{amsmath,amssymb,mathtools}
\usepackage[scr=rsfs,cal=euler]{mathalfa}
\usepackage{xfrac}
\usepackage[unicode,hidelinks,bookmarks=false]{hyperref}
\newcommand{\ie}{i.e.\ }
\newcommand{\cf}{cf.\ }
\newcommand{\resp}{respectively\ }
\newcommand{\Subsection}{\subsection}
\providecommand{\nobreakdash}{\nobreak\hbox{-}}
\setlength{\emergencystretch}{2em}
\multlinegap0pt
\usepackage{bm}
\usepackage{bbm}
\usepackage{dsfont}
\usepackage{xspace}
\usepackage{cancel}
\usepackage{mathtools}
\usepackage{amsthm}
\makeatletter
\@ifundefined{thm}{\newtheorem{thm}{Thm}}{}
\@ifundefined{lem}{\newtheorem{lem}[thm]{Lemma}}{}
\makeatother
\newcommand{\de}{\partial}
\title[Symmetry results for a nonlocal nonlinear Poincar\'e-Wirting]{Symmetry results for a nonlocal nonlinear Poincar\'e-Wirtinger inequality}
\author{Gianpaolo Piscitelli}
\date{Source: arXiv:2404.15486v3 (CC BY 4.0)}
\begin{document}
\maketitle

\noindent\emph{Source and licence.} Symmetry results for a nonlocal nonlinear Poincar\'e-Wirtinger inequality. Version arXiv:2404.15486v3 (CC BY 4.0), distributed under the Creative Commons Attribution 4.0 International licence, as recorded in the arXiv licence metadata for this exact version and shown on the abstract page https://arxiv.org/abs/2404.15486v3. Full rights notice: licenses/arxiv-2404.15486-CC-BY-4.0.txt.\par\smallskip

\noindent\textit{Editorial scope.} Counted unit: the Lem whose source label is \texttt{prop-monotonia-r} in the article (source0.tex lines 428--435, source label \texttt{prop-monotonia-r}), delivered together with the article's own complete original proof of it (source0.tex lines 436--561, one \texttt{proof} environment of 126 lines). No mathematics is written, repaired or re-derived: statement and proof are the article's own text line for line. Required context retained uncounted: Rdef (lines 416--418). Every definition, display and label of those lines is retained unchanged. Omitted from this package and not redistributed: 1--415, 419--427, 562--1067. Nothing inside a retained excerpt is omitted or altered: every retained body is delivered exactly as the article's lines are written, so no omission receipt of any kind accompanies this package. Citations: the retained excerpts carry no citation command at all, so no bibliography entry of the article is transcribed and no reference is redistributed. Reference adaptation: none was needed and none was made; every reference inside the delivered unit resolves inside it, so no unresolved reference or citation is produced. Printed slips kept literal: no printed word, symbol, hypothesis or number of the counted statement or of its proof was altered. Layout and class: the article's own class is \texttt{amsart} with packages including fontenc, inputenc, babel, graphicx, verbatim, color, picture, amsmath, amsthm, amsfonts, amssymb, mathtools, bigints, cleveref; the wrapper is the lane's pinned \texttt{amsart} editorial wrapper at 10pt on A4, which restates the theorem-like environments the retained text needs and adds the article's own macro definitions that the retained text needs (\texttt{\textbackslash de}), with extra wrapper packages (bm, bbm, dsfont, xspace, cancel, mathtools). The delivered numbering is the unit's own. Assets: no retained span contains a figure environment, a graphics-inclusion command, a PDF-page inclusion, a table or a TikZ picture; no image file is copied into this package, the package declares no asset dependency and no image file is redistributed with it. Licence: version arXiv:2404.15486v3 of this article is distributed under the Creative Commons Attribution 4.0 International licence, as recorded in the arXiv licence metadata for this exact version and shown on the abstract page https://arxiv.org/abs/2404.15486v3. A full rights notice is delivered at licenses/arxiv-2404.15486-CC-BY-4.0.txt. Characters: every retained excerpt is pure ASCII, so the delivered engine, XeLaTeX with its default Latin Modern text font, reports no missing character.
\begin{equation}\label{Rdef}
R(m,q,r)=\frac{1-m^q}{1+m^{r-1}}.
\end{equation}

\begin{lem}
\label{prop-monotonia-r}
Let $p,q,r>1$ such that $\frac 12+\frac q2 +\frac q{2p}\le r \le q+\frac q p$.  For any fixed $y\in[0,1[$ and
\begin{itemize}
\item for any fixed $m\in[0,1[$, the function $h(m,p,q,r,y)$ is strictly increasing with respect to $r$;
\item for $m=1$, the function $h(1,p,q,r,y)$ is constant with respect to $r$. 
\end{itemize}
\end{lem}

\begin{proof} We divide the proof into three steps: in the first step we compute the expression of the derivative of $h$ with respect to $r$ for any $m\in]0,1[$ and $y\in]0,1[$; in the second step we study the sign of the aforementioned derivative; in the third step we analyze the cases excluded by the previous steps. From now on, for the sake of simplicity, we will set  $R=R(m,q,r)$.

{\it Step 1 (The derivative of $h$).} Let us start by considering the case when $m\in]0,1[$ and $y\in]0,1[$.
Differentiating $h$ with respect to $r$, we have
\begin{equation*}
\begin{split}
\de_{r}h (m,p,q,r,y)= &- \frac 1 p \frac {(1-y^{r-1})\de_{r} R + R\, y^{r-1}\log y}{\left[1-R(1- y^{r-1}) - y^q\right]^\frac{p+1}p} +\\[.2cm]
&- \frac m p \frac{
[-(1+m^{r-1}y^{r-1})\de_{r}R- R\,m^{r-1}y^{r-1}(\log m +\log y)]}{ \left[1-R(1+m^{r-1}y^{r-1})-m^{q}y^{q}\right]^\frac {p+1} p}.
\end{split}
\end{equation*}
%\begin{equation*}\begin{split}\de_{r}h = &- \frac {1}{pF_{I}^{p+1}} \big[-(1-y^{r-1})\de_{r} R + R\, y^{r-1}\log y\big]+\\[.2cm]&- \frac {n}{pF_{I\!I}^{p+1}} \big[-(1+n^{r-1}y^{r-1})\de_{r}R- R\,n^{r-1}y^{r-1}(\log n +\log y)\big],\end{split}\end{equation*}
Therefore, in order to compute the derivative of $h$ with respect to $r$, we need to differentiate $R$ (defined in \eqref{Rdef}) with respect to $r$. We have
\[
\de_{r}R= -\frac{1-m^{q}}{(1+m^{r-1})^{2}}m^{r-1}\log m
\]
and hence
\begin{equation}
\begin{split}
\label{derivata_h}
\de_{r}h (m,p,q,r,y)= &-\frac1p\frac{1-m^{q}}{(1+m^{r-1})^{2}}\bigg\{\frac{ (1-y^{r-1})m^{r-1}\log m +y^{r-1}(1+m^{r-1})\log y}{\left[1-R(1- y^{r-1}) - y^q\right]^\frac{p+1}p}+ \\[.3cm]
& +m^r \frac{(1+m^{r-1}y^{r-1})\log m-(1+m^{r-1})y^{r-1}(\log m +\log y)}{ \left[1-R(1+m^{r-1}y^{r-1})-m^{q}y^{q}\right]^\frac {p+1} p}\bigg\}.
\end{split}
\end{equation}
{\it Step 2 (The monotonicity of $h$).}
It is easily seen that the numerator of the first ratio, in the curly brackets of \eqref{derivata_h}, is negative. If the numerator of the second ratio is also negative, we get the desired monotonicity. %\begin{equation}\label{segno_der_h}\partial _r h(n,p,q,r,y)>0.\end{equation}
Otherwise, if this second numerator is positive, let us observe that
\begin{equation*}
m^q(1-R(1- y^{r-1}) - y^q)\leq [1-R(1+m^{r-1}y^{r-1})-m^{q}y^{q}],
\end{equation*}
%\begin{equation*}\label{firint}F_{I}(n,p,q,r,y):=\displaystyle \left[1-R(1- y^{r-1}) - y^q\right]^\frac 1 p \le \left[1-y^q\right]^\frac 1p,\end{equation*}and \begin{equation*}\label{secint}F_{I\!I}(n,p,q,r,y):=\displaystyle \left[1-R(1+n^{r-1}y^{r-1})-n^{q}y^{q}\right]^\frac 1 p\ge n^q\left(1-y^{q}\right)^\frac 1p,\end{equation*}
%Resuming, we have\begin{multline*}\de_{r}h = \frac1p\frac{1-n^{q}}{(1+n^{r-1})^{2}}\bigg\{ \overbrace{\bigg[ -(1-y^{r-1})n^{r-1}\log n -y^{r-1}(1+n^{r-1})\log y\bigg]}^{h_{1}(n,r,y)}\frac{1}{F_{I}^{p+1}}+ \\[.3cm]+ \underbrace{\bigg[-(1+m^{r-1}y^{r-1})\log n+(1+n^{r-1})y^{r-1}(\log n +\log y)\bigg]}_{h_{2}(n,r,y)}\frac {n^{r}}{F_{I\!I}^{p+1}}\bigg\}.\end{multline*}
%By \eqref{firint} and \eqref{secint} 
that implies:
\begin{multline}\label{derivata_hge}
\de_{r}h(m,p,q,r,y) \ge -\frac1p\frac{1-m^{q}}{(1+m^{r-1})^{2}}\bigg\{\frac{ (1-y^{r-1})m^{r-1}\log m +y^{r-1}(1+m^{r-1})\log y}{m^{\frac {-q(p+1)}p}  \left[1-R(1+m^{r-1}y^{r-1})-m^{q}y^{q}\right]^\frac {p+1} p }+ \\[.3cm]
+m^r \frac{(1+m^{r-1}y^{r-1})\log m-(1+m^{r-1})y^{r-1}(\log m +\log y)}{ \left[1-R(1+m^{r-1}y^{r-1})-m^{q}y^{q}\right]^\frac {p+1} p}\bigg\}.
\end{multline}
%\begin{equation*}\begin{split}\de_{r}h \ge \frac1p\frac{1-n^{q}}{(1+n^{r-1})^{2}}\bigg\{ &\bigg[ -(1-y^{r-1})n^{r-1}\log n -y^{r-1}(1+n^{r-1})\log y\bigg]\frac{1}{(1-y^{p})^{\frac {p+1}{p}}}+ \\[.2cm]+& \bigg[(y^{r-1}-1)\log n+(1+n^{r-1})y^{r-1}\log y\bigg]\frac {n^{r-1-q}}{(1-y^{p})^{\frac{p+1}{p}}}\bigg\}.\end{split}\end{equation*}
Hence, by setting
\begin{multline*}
g(m,p,q,r,y):=\bigg[ -(1-y^{r-1})m^{r-1}\log m -y^{r-1}(1+m^{r-1})\log y\bigg]+\\
+ \bigg[(y^{r-1}-1)\log m+(1+m^{r-1})y^{r-1}\log y\bigg]m^{r-\frac{q(p+1)}p},
\end{multline*}
we have that inequality \eqref{derivata_hge} can be written as
\begin{equation}\label{derivative_h}
\de_{r}h (m,p,q,r,y) \ge \frac1p\frac{1-m^{q}}{(1+m^{r-1})^{2}m^{\frac {q(p+1)}p}} g(m,p,q,r,y).
\end{equation}
To prove the positivity of $\partial_r h$, we will show that
\begin{equation}
\label{segno_der_g}
g(m,p,q,r,y)>0,
\end{equation}
by proving that $g$ is decreasing with respect to $y$ in the interval $]0,1[$, since $g(m,p,q,r,1)=0$. By differentiating $g$ with respect to $y$, we obtain
\begin{equation*}
\begin{split}
\de_{y}g (m,p,q,r,y)& = \bigg[(r-1)y^{r-2}m^{r-1}\log m-(r-1)y^{r-2}(1+m^{r-1})\log y-y^{r-2}(1+m^{r})\bigg]\\
&\ \quad +\bigg[(r-1)y^{r-2}\log m+(1+m^{r-1})((r-1)y^{r-2}\log y+y^{r-2})\bigg]m^{r-\frac{q(p+1)}p}\\
& =y^{r-2}\bigg[(r-1) (m^{r-1} +m^{r-\frac{q(p+1)}p})\log m+(r-1)(1+m^{r-1})\left(m^{r-\frac{q(p+1)}p}-1\right)\log y  \\
&\ \quad +(1+m^{r-1})\left(m^{r-\frac{q(p+1)}p}-1\right)\bigg].
\end{split}
\end{equation*}
Since $r-\frac{q(p+1)}p\le  0$ by assumptions, this derivative is negative if and only if
%\begin{equation*}(1+n^{r-1})(n^{r-1-q}-1)((r-1)\log z+1) < (1-r)(n^{r-1}+n^{r-1-q})\log n.\end{equation*}The above inequality is true, as we will show that 
\begin{equation}
\label{ineqlog}
\log y < -\frac{\left(m^{r-1}+m^{r-\frac{q(p+1)}p}\right)\log m}{\left(1+m^{r-1}\right)\left(m^{r-\frac{q(p+1)}p}-1\right)}-\frac1{r-1}. %=:-\frac1r+\ell(n,q,r). 
\end{equation}
Since the left-hand term is negative, then if the right-hand side of \eqref{ineqlog} is nonnegative, then the inequality \eqref{ineqlog} holds.
%{\bf Claim 2.} {\em For any $q\in [p,+\infty[$, $r\in \left[\frac q2+1,q+1\right]$ and $n\in]0,1[$, $\ell(n,q,r)> \frac1{r-1}$.} We will show that\[\ell(n,q,r) >  \frac 1{r-1}.\]We have\[\ell(n,q,r)= \frac{(n^{r-1}+n^{r-1-q})\log n }{(1+n^{r-1})(1-n^{r-1-q})}> \frac 1{r-1}\]
To this aim, we will equivalently show that
\begin{multline}\label{segno_f}
f(m,p,q,r):=-\left(m^{r-1}+m^{r-\frac{q(p+1)}p}\right)\log m - \frac1{r-1}\left(1+m^{r-1}\right)\left(m^{r-\frac{q(p+1)}p}-1\right) > 0.
\end{multline}
We have
\begin{equation*}
\begin{split}
f(m,p,q,r)&=\left(m^{r-1}+m^{r-\frac{q(p+1)}p}\right)\log\frac{1}{m} +\frac1{r-1} \left( 1+m^{r-1}-m^{r-\frac{q(p+1)}p}-m^{2r-1-\frac{q(p+1)}p}\right)\\
&= m^{r-1}\left(\log\frac 1m+\frac1{r-1} \right)+ m^{r-\frac{q(p+1)}p}\left(\log\frac1m-\frac1{r-1} \right)+\frac1{r-1} \left(1-m^{2r-1-\frac{q(p+1)}p}\right)\\
&\ge m^{r-1}\left(\log\frac 1m+\frac1{r-1}\right)+ m^{r-\frac{q(p+1)}p}\left(\log\frac1m-\frac1{r-1}\right)\\
&=m^{r-\frac{q(p+1)}p}\left(m^{\frac{q(p+1)}p-1}\left(\log\frac 1m+\frac1{r-1}\right)+\log\frac 1m-\frac1{r-1}\right).
\end{split}
\end{equation*}
because $2r-1-\frac{q(p+1)}p\ge0$.
Hence, the positivity of $f$ follows if
\begin{equation}
\label{segno_e}
e(m,p,q,r):=m^{\frac{q(p+1)}p-1}\left(\log\frac 1m+\frac1{r-1}\right)+\log\frac 1m-\frac1{r-1}>0.
\end{equation}
Since $e(1,p,q,r)=0$, to prove \eqref{segno_e}, we show that $e$ is decreasing with respect to $m$; indeed, we have
\[
\partial_m e (n,p,q,r)=m^{\frac{q(p+1)}p-2}\left(\log\left( \frac 1{m^{\frac{q(p+1)}p-1}} \right)+\frac {\frac{q(p+1)}p-1}{r-1} -1-\frac1{m^{\frac{q(p+1)}p-1}} \right)
\]
that is negative since $\log z < z-1$ when $z>1$, $m^{\frac{q(p+1)}p-1}<1$ and $r\ge\frac 12+\frac q2 +\frac q{2p}$, by assumptions.

Hence \eqref{segno_e}, \eqref{segno_f}, \eqref{ineqlog} and \eqref{segno_der_g} are satisfied, and recalling the behavior of $h$ from \eqref{derivative_h}, this implies that
\[
\de_{r}h (m,p,q,r,y)\ge \frac1p\frac{1-m^{q}}{(1+m^{r-1})^{2}} g(m,p,q,r,y)> \frac1p\frac{1-m^{q}}{(1+m^{r-1})^{2}} 
g(m,p,q,r,1)=0.
\]
when $m\in]0,1[$ and $y\in]0,1[$.
%Finally, we show that if $m<1$ then $\de_{r}h(r,m,y)>0$ for $m\in ]0,1[$ and $y\in]0,1[$. This follows from the third claim below.
%
%{\bf Claim 3.} {\em If $0<m<1$ and $1<r\le 2$, then $g(r,m,\cdot$) is strictly decreasing for $y\in]0,1[$}.
%
%Taking a closer look to the inequality \eqref{ineqlog},
%
% we have that if $m<1$ the inequality \eqref{ineq0} is strict.

\textit{Step 3 (The trivial cases).} We observe that if $m=0$, then $R=1$ and 
\[h(0,p,q,r,y)=\frac{1}{ (y^{r-1}- y^q)^\frac 1 p},
\]
that is strictly increasing with respect to $r$.

Meanwhile, if $m=1$, then $R=0$ and
\[
h(1,p,q,r,y)=\frac{2}{ (1- y^q)^\frac 1 p}, 
\]
that is constant with respect to $r$. 

Finally, when $y=0$, we have
\[
h(m,p,q,r,0)=\frac{1+m}{1-R},
\]
that is strictly increasing with respect to $r$.
\end{proof}

\end{document}
